\section{Data Assets}
\label{sec:DataAssets}

Durga treats BPMN data objects as logical pipeline assets, not as Kafka
topics. Sequence flows still drive execution and Kafka still carries process
events, worker input/output messages, BPMN messages and BPMN signals. Data
objects describe the business artifacts that tasks read and write.

\subsection{Model}

A BPMN data object represents a named dataset or artifact, for example
\texttt{RawCustomerData}, \texttt{NormalizedCustomerData}, or
\texttt{CustomerGraphDelta}. A BPMN data store represents a physical source or
target such as S3, PostgreSQL, Neo4j, a local file store, or an API.

Data input and output associations connect executable activities to those
assets:

\begin{verbatim}
Transform Data
  input:  RawCustomerData
  output: NormalizedCustomerData

Enrich Data
  input:  FilteredCustomerData
  output: EnrichedCustomerData
  stores: S3Warehouse, Neo4jCustomerGraph, PostgresCustomerStore
\end{verbatim}

The scaffolder collects these declarations into \filepath{summary.json} and
the generated README. The declarations are metadata and contracts; they do not
change routing by themselves.

\subsection{Runtime handles}

Large datasets should normally travel by reference rather than as Kafka message
payload bytes. The generated runtime includes \texttt{DataHandle}:

\begin{verbatim}
DataHandle(
    name,
    uri,
    mediaType,
    schema,
    metadata)
\end{verbatim}

A process event payload can carry a serialized data handle such as an S3 URI,
media type, schema identifier, row count, checksum, or partition metadata. This
keeps lifecycle events small while preserving enough information for monitoring
and lineage.

\subsection{Extension properties}

Data objects and data stores may use Camunda extension properties for Durga
metadata:

\begin{verbatim}
schema    = schemas/enriched-customer-data.schema.json
mediaType = application/json
kind      = s3
uri       = s3://warehouse/customer/enriched/
\end{verbatim}

For data stores, \texttt{kind} identifies the adapter family. Current examples
include \texttt{s3}, \texttt{postgresql}, and \texttt{neo4j}. If \texttt{kind}
is omitted, later tooling may infer it from the URI scheme.

\subsection{Lineage direction}

The immediate implementation records data objects, data stores, and task
associations in generated metadata. The next runtime step is to promote this
metadata into lineage events, so task completions can report:

\begin{verbatim}
read:   [RawCustomerData]
wrote:  [EnrichedCustomerData]
stores: [S3Warehouse, Neo4jCustomerGraph]
\end{verbatim}

This keeps the process engine focused on orchestration while giving pipeline
operators the data-plane visibility needed for real workloads.
